\documentclass[12pt]{amsart}

\def\CHT{05} \def\SEC{00} \def\SEM{S} \def\YEAR{2019} \def\CLASS{MA311}

\newcommand{\BACKROOT}{../../../..}
\newcommand{\TEACHINGROOT}{\BACKROOT}
%\newcommand{\TEACHINGROOT}{/Users/rhcocker/Dropbox/Teaching}
\newcommand{\COURSEROOT}{\TEACHINGROOT/\SEM_\YEAR/\CLASS}
\newcommand{\SECTIONROOT}{\COURSEROOT/Chapter_\CHT/Section_\CHT_\SEC}
\newcommand{\PATHTOHEADERS}{\TEACHINGROOT/Headers}
\newcommand{\PATHTOSLIDES}{\SECTIONROOT/Slides}
\newcommand{\PATHTOFIGS}{\COURSEROOT/Figures}

\usepackage{mathrsfs}
\newcommand{\LAP}{\mathscr{L}}

\input{\PATHTOHEADERS/TestQuizPackageHeader.tex}

\renewcommand{\headrulewidth}{0pt} % remove lines as well

\numberwithin{equation}{section}

\DeclareGraphicsRule{.tif}{png}{.png}{`convert #1 `dirname #1`/`basename #1 .tif`.png}

\textheight = 10.8 in
\topmargin = -1.0in
\textwidth = 7.6in
%\oddsidemargin = -0.7in
%\evensidemargin = -0.7in

\newcounter{ProbNum}
\setcounter{ProbNum}{0}
\def\NEWPROB{\stepcounter{ProbNum} {\noindent \arabic{ProbNum}. \,\,}}

\begin{document}

\input{\PATHTOHEADERS/TestDefinitionHeader.tex}

\pagestyle{empty}

\noindent\rule{7.6in}{0.4pt}

{\centering\sc 
\noindent Differential Equations\\
 Concept Review for Laplace Transforms \\
}
\noindent\rule{7.6in}{0.4pt}

\vspace{0.1cm}

\UL{Review: Limits}

\NEWPROB Evaluate each of the following limits.
\begin{itemize}	
	\item  $\D \lim_{b\to \infty} \frac{1}{12} = $\\ \vspace{0.85cm}
	
	\item  $\D \lim_{b\to \infty} \frac{1}{s} = $ \\ \vspace{0.85cm}
	
	\item $\D \lim_{b\to \infty} e^{2b} = $\\ \vspace{0.85cm}
	
	\item $\D \lim_{b\to \infty} e^{0.1b} = $\\ \vspace{0.85cm}
	
	\item $\D \lim_{b\to \infty} e^{-2b} = $\\ \vspace{0.85cm}
\end{itemize}
	
\vspace{0.7cm}

\NEWPROB For what values of $a$ does $\D \lim_{b \to \infty}e^{ab}$ 
	converge?\\ 

\newpage
\thispagestyle{empty}

\NEWPROB For each of the following, give the condition (if necessary) when the limit converges and find the limit.
\begin{itemize}
	\item  $\D \lim_{b\to \infty} \left[ \frac{1}{2}e^{2b} - \frac{1}{2} \right] = $\\ \vspace{1.75cm}
	
	\item  $\D \lim_{b\to \infty} \left[ -\frac{1}{-3}e^{-3b} + \frac{1}{-3} \right] = $\\ \vspace{1.75cm}
	
	\item  $\D \lim_{b\to \infty} \left[ -\frac{1}{s}e^{-sb} + \frac{1}{s} \right] = $ \\ \vspace{1.75cm}
	
	\item  $\D \lim_{b\to \infty} \left[ -\frac{1}{s}e^{-sb} + \frac{1}{s} \right] = $ \\ \vspace{1.75cm}
	
	\item  $\D \lim_{b\to \infty} \left[ \frac{1}{s^2}e^{-sb} \right] = $  \\ \vspace{1.75cm}
	
	\item  $\D  \lim_{b\to \infty} \left[ \frac{1}{s-3}e^{(3-s)b} + \frac{1}{3-s} \right] = $ \\ \vspace{1.75cm}
	
	\item  $\D  \lim_{b\to \infty} \left[ \frac{1}{s+7}e^{(-7-s)b} + \frac{1}{-7-s} \right] = $ \\ \vspace{1.75cm}
	
	\item  $\D \lim_{b\to \infty} \left[ \frac{1}{s-a}e^{(a-s)b} + \frac{1}{a-s} \right] = $ \\ \vspace{1.75cm}
\end{itemize}

\newpage
\thispagestyle{empty}

%%%%%%%%%%%%%%%%%%%%%%%%%%%
\UL{Review: L'Hospital's Rule} 

Recall: This rule can only be applied when the limit has an indeterminate form like $\frac{0}{0}$ or $\frac{\infty}{\infty}$.

\NEWPROB Evaluate each of the following limits.
\begin{itemize}
	\item  $\D \lim_{b\to \infty} \frac{b}{e^b}$ \\ \vspace{3.55cm}
	\item  $\D \lim_{b\to \infty} be^{-3b}$ \\ \vspace{3.55cm}
	\item  $\D \lim_{b\to \infty} \left[ -\frac{1}{5}be^{-5b} - \frac{1}{25}e^{-5b} - \frac{1}{25} \right]$ \\ \vspace{3.55cm}
	\item  $\D \lim_{b\to \infty} \left[ -\frac{1}{s}be^{-sb} - \frac{1}{s^2}e^{-sb} - \frac{1}{s^2} \right]$ \\ \vspace{3.55cm}
\end{itemize}
\vspace{0.5cm}

\newpage
\thispagestyle{empty}

%%%%%%%%%%%%%%%%%%%%%%%%%%%
\UL{Review: Improper Integrals}

\NEWPROB Evaluate the following integrals using a limit.
\begin{itemize}
	\item  $\D \int_0^{\infty}e^{-3t}dt = $\\ \vspace{8.25cm}
	\item  $\D \int_0^{\infty}e^{-st}dt = $\\ \vspace{4.25cm}
\end{itemize}

\newpage
\thispagestyle{empty}

%\UL{Laplace Transform}
%
%\noindent  The Laplace Transform of function $f(t)$ is a new function $F(s)$ defined by:
%	\[ F(s) = \LAP \{ f(t) \} = \int_0^{\infty}e^{-st}f(t) dt, \]
%	provided this integral is finite.\\
%	
%\NEWPROB Use the definition of Laplace Transform (above) to find the Laplace Transform of $h(t) = 7t.$  
%	Include appropriate restrictions on $s$ to guarantee the convergence of the integral.\\
%
%\NEWPROB Use the definition of Laplace Transform (above) to find the Laplace Transform of $w(t) = e^{kt},$ 
%	where $k$ is some constant.  Include appropriate restrictions on $s$ to guarantee the
%	convergence of the integral.\\
%	
%\NEWPROB Use the definition of Laplace Transform (above) to find the Laplace Transform of $y(t) = 7t + e^{5t}.$
%		Include appropriate restrictions on $s$ to guarantee the convergence of the integral.\\
%
%\NEWPROB Use the definition of Laplace Transform (above) to find the Laplace Transform of  the function
%	\[d(t) = \left\{ \begin{array}{ll}	0,	& t<0\\
%							7,	& 0 \le t < 3\\
%							0,	& t \ge 3\\
%							\end{array} \right.\]  %= 7U(t) - 7U(t-3)
%	Include appropriate restrictions on $s$ to guarantee the convergence of the integral.\\
%	
%%\item    $\D H(s) = \frac{7}{s^2},$ $s>0$\\
%%
%%\item    $\D W(s) = \frac{1}{s-k},$  $s > k$	\\
%%
%%\item   $\D Y(s) = \frac{2}{s^2} + \frac{1}{s-5},$  $s>5$
%
%%\item  $\D D(s) = \frac{7}{s} - \frac{7e^{-3s}}{s}$
%
%\newpage
%\thispagestyle{empty}
%
%\NEWPROB Use the table and properties provided to determine the Laplace transform of each function.
%Include resulting restrictions on $s.$
%%-------NUM LIST---------
%\begin{enumerate}
%	\item  $\D f(t) = e^{-2t}\sin(2t) + t^2 e^{3t}$\\ 
%	\item  $\D g(t) = t^3 e^{-9t}$\\
%	\item  $\D h(t) = e^{2t} - t^3 + t^2 - \sin (5t)$\\
%	\item  $\D q(t) = 8t\cos(6t) + e^{3t}\sin(4t)$\\
%	\item  $\D f(t) = t^2 \sin(3t)$\\
%\end{enumerate}
%		
%\newpage
%\thispagestyle{empty}

\end{document}